% GENERATED FILE. Edit canonical lessons, then run npm run generate:books.
% canonical-source-hash: fnv1a64:7cefb8ad045eb9dc
% canonical-lessons: PT-C425-albergue, PT-C425-pousada, PT-C425-estalagem, PT-C425-residencial, PT-C425-campismo

\chapter{Onde dormir}
\label{ch:pt-onde-dormir}
\hlchaptermodality{425}

\begin{chapteropening}
\textbf{By the end of this chapter:} \emph{I can choose where to sleep: a hostel, a country inn, a roadside inn, a guesthouse or a campsite.}

Complete the last lesson of chapter 425: I can choose where to sleep: a hostel, a country inn, a roadside inn, a guesthouse or a campsite.
\end{chapteropening}

\section[o albergue]{o albergue — the hostel}

\label{lesson:PT-C425-albergue}



\begin{quote}
Before the new one: say the Portuguese for the flight attendant, then the Portuguese for the crew.
\end{quote}

\subsection*{You'll want to know: o albergue}
\textbf{o albergue} — \textquotedblleft{}the hostel\textquotedblright{}.

\begin{quote}
Ficámos num albergue para jovens.
\end{quote}

\begin{quote}
\emph{We stayed in a hostel for young people.}
\end{quote}

Pilgrims on the long walks to the north sleep in \emph{albergues}.

\begin{cousinweb}
\emph{albergue} came through Occitan from Gothic \emph{haribairgō} 'shelter for an army' (EtymDB).
\end{cousinweb}

\begin{grammarlens}[title={What you've built}]
One more word for telling, explaining and describing at B1.
\end{grammarlens}

\subsection*{Your turn}
\begin{itemize}
\raggedright
  \item \emph{Say it:} \emph{o albergue}
  \item \emph{Say it:} \emph{o albergue}, once more
  \item \emph{Say it:} a sentence of your own with \emph{o albergue}
  \item \emph{Recall:} say the Portuguese for the flight attendant, then the Portuguese for the crew, then say \emph{o albergue} again
\end{itemize}

\begin{tcolorbox}[breakable,colback=teal!4,colframe=teal!35!black,title={Before you move on}]
What does o albergue mean? (\textquotedblleft{}The hostel\textquotedblright{}.) Say it once more.
\end{tcolorbox}
\section[a pousada]{a pousada — the country inn; the hotel in a historic building}

\label{lesson:PT-C425-pousada}



\begin{quote}
Before the new one: say the Portuguese for the crew, then the Portuguese for the hostel.
\end{quote}

\subsection*{You'll want to know: a pousada}
\textbf{a pousada} — \textquotedblleft{}the country inn; the hotel in a historic building\textquotedblright{}.

\begin{quote}
Dormimos numa pousada num castelo antigo.
\end{quote}

\begin{quote}
\emph{We slept in an inn in an old castle.}
\end{quote}

\emph{A pousada de juventude} — the youth hostel.

\begin{cousinweb}
\emph{pousada} is built on \emph{pousar} 'to set down, to rest', which is inherited from Latin \emph{pausāre} 'to halt, to rest' (EtymDB).
\end{cousinweb}

\begin{grammarlens}[title={What you've built}]
One more word for telling, explaining and describing at B1.
\end{grammarlens}

\subsection*{Your turn}
\begin{itemize}
\raggedright
  \item \emph{Say it:} \emph{a pousada}
  \item \emph{Say it:} \emph{a pousada}, once more
  \item \emph{Say it:} a sentence of your own with \emph{a pousada}
  \item \emph{Recall:} say the Portuguese for the crew, then the Portuguese for the hostel, then say \emph{a pousada} again
\end{itemize}

\begin{tcolorbox}[breakable,colback=teal!4,colframe=teal!35!black,title={Before you move on}]
What does a pousada mean? (\textquotedblleft{}The country inn; the hotel in a historic building\textquotedblright{}.) Say it once more.
\end{tcolorbox}
\section[a estalagem]{a estalagem — the roadside inn, the inn}

\label{lesson:PT-C425-estalagem}



\begin{quote}
Before the new one: say the Portuguese for the hostel, then the Portuguese for the country inn.
\end{quote}

\subsection*{You'll want to know: a estalagem}
\textbf{a estalagem} — \textquotedblleft{}the roadside inn, the inn\textquotedblright{}.

\begin{quote}
A estalagem fica à beira da estrada.
\end{quote}

\begin{quote}
\emph{The inn is by the side of the road.}
\end{quote}

An old word; today often the name of a small hotel.

\begin{cousinweb}
\emph{estalagem} came from Occitan \emph{ostalatge} 'lodging', built on \emph{ostal} 'house', from Late Latin \emph{hospitāle} 'guest house' (EtymDB).
\end{cousinweb}

\begin{grammarlens}[title={What you've built}]
One more word for telling, explaining and describing at B1.
\end{grammarlens}

\subsection*{Your turn}
\begin{itemize}
\raggedright
  \item \emph{Say it:} \emph{a estalagem}
  \item \emph{Say it:} \emph{a estalagem}, once more
  \item \emph{Say it:} a sentence of your own with \emph{a estalagem}
  \item \emph{Recall:} say the Portuguese for the hostel, then the Portuguese for the country inn, then say \emph{a estalagem} again
\end{itemize}

\begin{tcolorbox}[breakable,colback=teal!4,colframe=teal!35!black,title={Before you move on}]
What does a estalagem mean? (\textquotedblleft{}The roadside inn, the inn\textquotedblright{}.) Say it once more.
\end{tcolorbox}
\section[a residencial]{a residencial — the guesthouse}

\label{lesson:PT-C425-residencial}



\begin{quote}
Before the new one: say the Portuguese for the country inn, then the Portuguese for the roadside inn.
\end{quote}

\subsection*{You'll want to know: a residencial}
\textbf{a residencial} — \textquotedblleft{}the guesthouse\textquotedblright{}.

\begin{quote}
A residencial não é cara e fica perto da estação.
\end{quote}

\begin{quote}
\emph{The guesthouse is not expensive and is close to the station.}
\end{quote}

Also an adjective: \emph{um bairro residencial}.

\begin{cousinweb}
\emph{residencial} is built on \emph{residência} 'residence', a learned word that goes back to Latin \emph{residēre} 'to stay behind, to settle', on \emph{sedēre} 'to sit' (Lewis \& Short).
\end{cousinweb}

\begin{grammarlens}[title={What you've built}]
One more word for telling, explaining and describing at B1.
\end{grammarlens}

\subsection*{Your turn}
\begin{itemize}
\raggedright
  \item \emph{Say it:} \emph{a residencial}
  \item \emph{Say it:} \emph{a residencial}, once more
  \item \emph{Say it:} a sentence of your own with \emph{a residencial}
  \item \emph{Recall:} say the Portuguese for the country inn, then the Portuguese for the roadside inn, then say \emph{a residencial} again
\end{itemize}

\begin{tcolorbox}[breakable,colback=teal!4,colframe=teal!35!black,title={Before you move on}]
What does a residencial mean? (\textquotedblleft{}The guesthouse\textquotedblright{}.) Say it once more.
\end{tcolorbox}
\section[o campismo]{o campismo — camping}

\label{lesson:PT-C425-campismo}



\begin{quote}
Before the new one: say the Portuguese for the roadside inn, then the Portuguese for the guesthouse.
\end{quote}

\subsection*{You'll want to know: o campismo}
\textbf{o campismo} — \textquotedblleft{}camping\textquotedblright{}.

\begin{quote}
O parque de campismo tem duche e piscina.
\end{quote}

\begin{quote}
\emph{The campsite has showers and a pool.}
\end{quote}

\emph{O parque de campismo} — the campsite. In Brazil: \emph{o camping}.

\begin{cousinweb}
\emph{campismo} is built on \emph{campo} 'field' with \emph{-ismo}; \emph{campo} is inherited from Latin \emph{campus} 'field, open plain' (EtymDB).
\end{cousinweb}

\begin{grammarlens}[title={What you've built}]
One more word for telling, explaining and describing at B1.
\end{grammarlens}

\subsection*{Your turn}
\begin{itemize}
\raggedright
  \item \emph{Say it:} \emph{o campismo}
  \item \emph{Say it:} \emph{o campismo}, once more
  \item \emph{Say it:} a sentence of your own with \emph{o campismo}
  \item \emph{Recall:} say the Portuguese for the roadside inn, then the Portuguese for the guesthouse, then say \emph{o campismo} again
\end{itemize}

\begin{tcolorbox}[breakable,colback=teal!4,colframe=teal!35!black,title={Before you move on}]
What does o campismo mean? (\textquotedblleft{}Camping\textquotedblright{}.) Say it once more.
\end{tcolorbox}
